% Generado desde docs/INFORME.md. Editar libremente.

\section{6. Rutas, secretos, identidades y operaciones GraphQL}

Cada petición que llega de afuera debe terminar en el servicio correcto, y cada servicio necesita una forma de demostrar quién es. Aquí está el mapa completo: qué ruta pública va a qué servicio, con qué control y con qué credencial; qué secretos existen, quién los usa y cómo se cambian; y qué operaciones ofrece GraphQL. Las credenciales se nombran por su ubicación en Vault, nunca por su valor.


\subsection{Cómo se piensan las llaves}

La imagen que sirve aquí es la de un edificio con muchas puertas. Cada servicio tiene su propia llave y esa llave no abre las demás: si alguien consigue la de pagos, solo entra a pagos. El gateway es el único que lleva varias llaves, porque su trabajo es llamar a todos, y por eso es la pieza más vigilada. Los tres tipos de identidad que aparecen en esta sección son distintos entre sí: la \textbf{persona} se identifica con su cookie de sesión, un \textbf{servicio} con su credencial interna, y la \textbf{pasarela} con la firma de sus avisos.


Las llaves no están escritas en el código ni en el repositorio: viven en Vault, y \texttt{vault-\allowbreak{}init}, el contenedor que corre una sola vez al arrancar, las crea al azar. Cada servicio recibe un token que solo le permite leer sus propias rutas, y cada uno tiene su propio usuario de base de datos. Una limitación que está declarada entre las brechas de la sección 0: en Docker Compose los tokens de todos los servicios se escriben en un mismo volumen de solo lectura, en vez de montarse cada uno solo en su contenedor; lo que protege es que cada token solo puede leer sus propias rutas de Vault. En producción lo resolvería el orquestador (Docker secrets o Kubernetes). Si Vault deja de responder, el sistema se cierra en vez de abrirse: después de unos minutos de gracia, los servicios responden 503 y no dejan pasar a nadie.


\subsection{Matriz de rutas}

Cada fila es una operación del módulo 1 (mismo ID). Lo que no está en esta tabla no existe hacia afuera (404). Las rutas \texttt{/\allowbreak{}internos/\allowbreak{}*} de los servicios nunca se publican. Se generó leyendo \texttt{gateway/\allowbreak{}src/\allowbreak{}rutas.\allowbreak{}js}, que tiene 36 filas (S30 a S32 comparten la de \texttt{/\allowbreak{}graphql}); a ellas se suman 6 operaciones que resuelve el propio gateway (S33 a S37 y S44). Controles que se aplican a \textbf{todas} las rutas, además de los de la columna: HTTPS con HSTS, límite general por IP, origen permitido en los verbos que cambian datos (CSRF), cuerpo de hasta 100~KB, plazo de 8~s (504 si se vence, 502 si el servicio no responde), errores con la misma forma y \texttt{X-\allowbreak{}Request-\allowbreak{}Id} en los logs. Antes de reenviar, el gateway descarta \texttt{X-\allowbreak{}Usuario-\allowbreak{}*}, \texttt{X-\allowbreak{}Credencial-\allowbreak{}Interna}, \texttt{Authorization} y la cookie, y agrega solo la identidad verificada y la credencial del destino.


\begin{landscape}
\vspace{0.6em}
{\scriptsize
\setlength{\LTpre}{0pt}\setlength{\LTpost}{0pt}
\setlength{\tabcolsep}{3pt}\renewcommand{\arraystretch}{1.15}
\begin{longtable}{|>{\raggedright\arraybackslash}p{\dimexpr 0.0455\linewidth-2\tabcolsep\relax}|>{\raggedright\arraybackslash}p{\dimexpr 0.0711\linewidth-2\tabcolsep\relax}|>{\raggedright\arraybackslash}p{\dimexpr 0.0775\linewidth-2\tabcolsep\relax}|>{\raggedright\arraybackslash}p{\dimexpr 0.1278\linewidth-2\tabcolsep\relax}|>{\raggedright\arraybackslash}p{\dimexpr 0.1077\linewidth-2\tabcolsep\relax}|>{\raggedright\arraybackslash}p{\dimexpr 0.1745\linewidth-2\tabcolsep\relax}|>{\raggedright\arraybackslash}p{\dimexpr 0.2008\linewidth-2\tabcolsep\relax}|>{\raggedright\arraybackslash}p{\dimexpr 0.1950\linewidth-2\tabcolsep\relax}|}
\hline
\textbf{ID} & \textbf{Gateway} & \textbf{Verbo} & \textbf{Ruta externa (\texttt{/\allowbreak{}api/\allowbreak{}v1})} & \textbf{Acceso y rol} & \textbf{Sistema y ruta destino} & \textbf{Controles propios} & \textbf{Referencia de credencial} \\ \hline
\endfirsthead
\hline
\textbf{ID} & \textbf{Gateway} & \textbf{Verbo} & \textbf{Ruta externa (\texttt{/\allowbreak{}api/\allowbreak{}v1})} & \textbf{Acceso y rol} & \textbf{Sistema y ruta destino} & \textbf{Controles propios} & \textbf{Referencia de credencial} \\ \hline
\endhead
S01 & gateway externo & GET & \texttt{/\allowbreak{}catalogo/\allowbreak{}items} & Público & inventario: \texttt{/\allowbreak{}items} & Filtros validados (400) & \texttt{internos/\allowbreak{}inventario} \\ \hline
S02 & gateway externo & GET & \texttt{/\allowbreak{}catalogo/\allowbreak{}items/\allowbreak{}:\allowbreak{}id} & Público & inventario: \texttt{/\allowbreak{}items/\allowbreak{}:\allowbreak{}id} & — & \texttt{internos/\allowbreak{}inventario} \\ \hline
S03 & gateway externo & GET & \texttt{/\allowbreak{}catalogo/\allowbreak{}configuracion} & Público & inventario: \texttt{/\allowbreak{}configuracion} & — & \texttt{internos/\allowbreak{}inventario} \\ \hline
S04 & gateway externo & POST & \texttt{/\allowbreak{}catalogo/\allowbreak{}items} & admin, dueño & inventario: \texttt{/\allowbreak{}items} & Validación de producto & \texttt{internos/\allowbreak{}inventario} \\ \hline
S05 & gateway externo & PUT & \texttt{/\allowbreak{}catalogo/\allowbreak{}items/\allowbreak{}:\allowbreak{}id} & admin, dueño & inventario: \texttt{/\allowbreak{}items/\allowbreak{}:\allowbreak{}id} & Validación de producto & \texttt{internos/\allowbreak{}inventario} \\ \hline
S06 & gateway externo & DELETE & \texttt{/\allowbreak{}catalogo/\allowbreak{}items/\allowbreak{}:\allowbreak{}id} & admin & inventario: \texttt{/\allowbreak{}items/\allowbreak{}:\allowbreak{}id} & Solo admin en gateway e inventario & \texttt{internos/\allowbreak{}inventario} \\ \hline
S07 & gateway externo & POST & \texttt{/\allowbreak{}clientes/\allowbreak{}registro} & Público & clientes: \texttt{/\allowbreak{}clientes/\allowbreak{}registro} & 20/min por IP; validación de ficha & \texttt{internos/\allowbreak{}clientes} \\ \hline
S08 & gateway externo & GET & \texttt{/\allowbreak{}clientes/\allowbreak{}yo} & cliente & clientes: \texttt{/\allowbreak{}clientes/\allowbreak{}yo} & Propiedad & \texttt{internos/\allowbreak{}clientes} \\ \hline
S09 & gateway externo & PUT & \texttt{/\allowbreak{}clientes/\allowbreak{}yo} & cliente & clientes: \texttt{/\allowbreak{}clientes/\allowbreak{}yo} & Propiedad & \texttt{internos/\allowbreak{}clientes} \\ \hline
S10 & gateway externo & GET & \texttt{/\allowbreak{}clientes} & admin & clientes: \texttt{/\allowbreak{}clientes} & — & \texttt{internos/\allowbreak{}clientes} \\ \hline
S11 & gateway externo & POST & \texttt{/\allowbreak{}clientes} & admin & clientes: \texttt{/\allowbreak{}clientes} & Clave temporal & \texttt{internos/\allowbreak{}clientes} \\ \hline
S12 & gateway externo & GET & \texttt{/\allowbreak{}clientes/\allowbreak{}:\allowbreak{}id} & admin & clientes: \texttt{/\allowbreak{}clientes/\allowbreak{}:\allowbreak{}id} & — & \texttt{internos/\allowbreak{}clientes} \\ \hline
S13 & gateway externo & PUT & \texttt{/\allowbreak{}clientes/\allowbreak{}:\allowbreak{}id} & admin & clientes: \texttt{/\allowbreak{}clientes/\allowbreak{}:\allowbreak{}id} & — & \texttt{internos/\allowbreak{}clientes} \\ \hline
S14 & gateway externo & PATCH & \texttt{/\allowbreak{}clientes/\allowbreak{}:\allowbreak{}id/\allowbreak{}estado} & admin & clientes: \texttt{/\allowbreak{}clientes/\allowbreak{}:\allowbreak{}id/\allowbreak{}estado} & Cierra sesiones al desactivar & \texttt{internos/\allowbreak{}clientes} \\ \hline
S15 & gateway externo & POST & \texttt{/\allowbreak{}pedidos} & cliente & pedidos: \texttt{/\allowbreak{}pedidos} & \texttt{Idempotency-\allowbreak{}Key}; precio de inventario & \texttt{internos/\allowbreak{}pedidos} \\ \hline
S16 & gateway externo & GET & \texttt{/\allowbreak{}pedidos} & Todos los roles & pedidos: \texttt{/\allowbreak{}pedidos} & Propiedad y filtro por rol & \texttt{internos/\allowbreak{}pedidos} \\ \hline
S17 & gateway externo & GET & \texttt{/\allowbreak{}pedidos/\allowbreak{}:\allowbreak{}id} & Todos los roles & pedidos: \texttt{/\allowbreak{}pedidos/\allowbreak{}:\allowbreak{}id} & Propiedad y filtro por rol & \texttt{internos/\allowbreak{}pedidos} \\ \hline
S18 & gateway externo & GET & \texttt{/\allowbreak{}pedidos/\allowbreak{}boleta/\allowbreak{}:\allowbreak{}numero} & cliente, admin, dueño & pedidos: \texttt{/\allowbreak{}pedidos/\allowbreak{}boleta/\allowbreak{}:\allowbreak{}numero} & Propiedad & \texttt{internos/\allowbreak{}pedidos} \\ \hline
S19 & gateway externo & GET & \texttt{/\allowbreak{}pedidos/\allowbreak{}configuracion} & Todos los roles & pedidos: \texttt{/\allowbreak{}pedidos-\allowbreak{}config} & — & \texttt{internos/\allowbreak{}pedidos} \\ \hline
S20 & gateway externo & POST & \texttt{/\allowbreak{}pedidos/\allowbreak{}:\allowbreak{}id/\allowbreak{}aceptar} & admin, cocina & pedidos: \texttt{/\allowbreak{}pedidos/\allowbreak{}:\allowbreak{}id/\allowbreak{}aceptar} & Estado atómico & \texttt{internos/\allowbreak{}pedidos} \\ \hline
S21 & gateway externo & POST & \texttt{/\allowbreak{}pedidos/\allowbreak{}:\allowbreak{}id/\allowbreak{}terminar} & admin, cocina & pedidos: \texttt{/\allowbreak{}pedidos/\allowbreak{}:\allowbreak{}id/\allowbreak{}terminar} & Estado atómico & \texttt{internos/\allowbreak{}pedidos} \\ \hline
S22 & gateway externo & POST & \texttt{/\allowbreak{}pedidos/\allowbreak{}:\allowbreak{}id/\allowbreak{}despachar} & admin, despacho & pedidos: \texttt{/\allowbreak{}pedidos/\allowbreak{}:\allowbreak{}id/\allowbreak{}despachar} & Estado atómico & \texttt{internos/\allowbreak{}pedidos} \\ \hline
S23 & gateway externo & POST & \texttt{/\allowbreak{}pedidos/\allowbreak{}:\allowbreak{}id/\allowbreak{}entregar} & admin, despacho & pedidos: \texttt{/\allowbreak{}pedidos/\allowbreak{}:\allowbreak{}id/\allowbreak{}entregar} & Estado atómico & \texttt{internos/\allowbreak{}pedidos} \\ \hline
S24 & gateway externo & POST & \texttt{/\allowbreak{}pedidos/\allowbreak{}:\allowbreak{}id/\allowbreak{}impresa} & admin, cocina, despacho & pedidos: \texttt{/\allowbreak{}pedidos/\allowbreak{}:\allowbreak{}id/\allowbreak{}impresa} & — & \texttt{internos/\allowbreak{}pedidos} \\ \hline
S25 & gateway externo & POST & \texttt{/\allowbreak{}pedidos/\allowbreak{}:\allowbreak{}id/\allowbreak{}anular} & cliente, admin & pedidos: \texttt{/\allowbreak{}pedidos/\allowbreak{}:\allowbreak{}id/\allowbreak{}anular} & Propiedad, estado y motivo & \texttt{internos/\allowbreak{}pedidos} \\ \hline
S26 & gateway externo & GET & \texttt{/\allowbreak{}reportes/\allowbreak{}ventas} & admin, dueño & pedidos: \texttt{/\allowbreak{}reportes/\allowbreak{}ventas} & Período $\leq$ 1 año & \texttt{internos/\allowbreak{}pedidos} \\ \hline
S27 & gateway externo & POST & \texttt{/\allowbreak{}pagos} & cliente & pagos: \texttt{/\allowbreak{}pagos} & \texttt{Idempotency-\allowbreak{}Key} obligatoria; propiedad vía pedidos & \texttt{internos/\allowbreak{}pagos} \\ \hline
S28 & gateway externo & GET & \texttt{/\allowbreak{}pagos} & cliente, admin, dueño & pagos: \texttt{/\allowbreak{}pagos} & Propiedad vía pedidos & \texttt{internos/\allowbreak{}pagos} \\ \hline
S29 & gateway externo & POST & \texttt{/\allowbreak{}pagos/\allowbreak{}notificacion} & Sin sesión (firma de la pasarela) & pagos: \texttt{/\allowbreak{}pagos/\allowbreak{}notificacion} & Firma HMAC; se reenvía sin identidad de usuario & \texttt{internos/\allowbreak{}pagos}; firma con \texttt{pasarela/\allowbreak{}comercio} (\texttt{secreto\_\allowbreak{}firma}) \\ \hline
S30 & gateway externo & POST & \texttt{/\allowbreak{}graphql} \texttt{resumenPedido} & Con sesión & GraphQL Gateway: \texttt{/\allowbreak{}graphql} & Profundidad $\leq$ 6, costo $\leq$ 1.000, paginación $\leq$ 50 & \texttt{internos/\allowbreak{}graphql} \\ \hline
S31 & gateway externo & POST & \texttt{/\allowbreak{}graphql} \texttt{misPedidos} & Con sesión & GraphQL Gateway: \texttt{/\allowbreak{}graphql} & Profundidad, costo y paginación & \texttt{internos/\allowbreak{}graphql} \\ \hline
S32 & gateway externo & POST & \texttt{/\allowbreak{}graphql} \texttt{reporteVentas} & Con sesión & GraphQL Gateway: \texttt{/\allowbreak{}graphql} & Profundidad, costo y paginación & \texttt{internos/\allowbreak{}graphql} \\ \hline
S33 & gateway externo & POST & \texttt{/\allowbreak{}auth/\allowbreak{}login} & Público & Auth Service: \texttt{/\allowbreak{}login} & 20/min por IP; bloqueo por correo & \texttt{internos/\allowbreak{}auth-\allowbreak{}gateway} (Basic) \\ \hline
S34 & gateway externo & POST & \texttt{/\allowbreak{}auth/\allowbreak{}logout} & Público & Auth Service: \texttt{/\allowbreak{}logout} & Borra la cookie & \texttt{internos/\allowbreak{}auth-\allowbreak{}gateway} (Basic) \\ \hline
S35 & gateway externo & GET & \texttt{/\allowbreak{}auth/\allowbreak{}sesion} & Público & gateway (no reenvía) & Respuesta del propio gateway & — \\ \hline
S36 & gateway externo & GET & \texttt{/\allowbreak{}auth/\allowbreak{}config} & Público & gateway (no reenvía) & Respuesta del propio gateway & — \\ \hline
S37 & gateway externo & PUT & \texttt{/\allowbreak{}auth/\allowbreak{}clave} & Con sesión & Auth Service: \texttt{/\allowbreak{}clave} & Cuenta tomada de la sesión & \texttt{internos/\allowbreak{}auth-\allowbreak{}gateway} (Basic) \\ \hline
S38 & gateway externo & POST & \texttt{/\allowbreak{}auth/\allowbreak{}verificar} & Público & Auth Service: \texttt{/\allowbreak{}verificacion} & 20/min por IP; respuesta uniforme & \texttt{internos/\allowbreak{}auth-\allowbreak{}gateway} (Basic) \\ \hline
S39 & gateway externo & POST & \texttt{/\allowbreak{}auth/\allowbreak{}reenviar} & Público & Auth Service: \texttt{/\allowbreak{}verificacion/\allowbreak{}reenviar} & 20/min por IP; respuesta uniforme & \texttt{internos/\allowbreak{}auth-\allowbreak{}gateway} (Basic) \\ \hline
S40 & gateway externo & GET & \texttt{/\allowbreak{}usuarios} & admin & Auth Service: \texttt{/\allowbreak{}usuarios} & — & \texttt{internos/\allowbreak{}auth-\allowbreak{}gateway} (Basic) \\ \hline
S41 & gateway externo & POST & \texttt{/\allowbreak{}usuarios} & admin & Auth Service: \texttt{/\allowbreak{}usuarios} & Clave fuerte & \texttt{internos/\allowbreak{}auth-\allowbreak{}gateway} (Basic) \\ \hline
S42 & gateway externo & PUT & \texttt{/\allowbreak{}usuarios/\allowbreak{}:\allowbreak{}id} & admin & Auth Service: \texttt{/\allowbreak{}usuarios/\allowbreak{}:\allowbreak{}id} & Último admin y propio perfil protegidos & \texttt{internos/\allowbreak{}auth-\allowbreak{}gateway} (Basic) \\ \hline
S43 & gateway externo & PATCH & \texttt{/\allowbreak{}usuarios/\allowbreak{}:\allowbreak{}id/\allowbreak{}estado} & admin & Auth Service: \texttt{/\allowbreak{}usuarios/\allowbreak{}:\allowbreak{}id/\allowbreak{}estado} & Último admin y propia cuenta protegidos & \texttt{internos/\allowbreak{}auth-\allowbreak{}gateway} (Basic) \\ \hline
S44 & gateway externo & GET & \texttt{/\allowbreak{}api/\allowbreak{}salud} & Público & gateway: \texttt{GET /\allowbreak{}salud} de cada servicio y de Vault & Sin datos personales; consulta cada \texttt{/\allowbreak{}salud} con 2~s de plazo & — \\ \hline
\end{longtable}
}
\end{landscape}

\begin{landscape}
Las referencias de credencial son rutas simbólicas en Vault (\texttt{secret/\allowbreak{}poke-\allowbreak{}fresh/\allowbreak{}.\allowbreak{}.\allowbreak{}.\allowbreak{}}), no valores. Los tres GraphQL (S30 a S32) comparten la ruta \texttt{POST /\allowbreak{}graphql}: el gateway solo exige sesión y \textbf{cada operación y cada campo se autoriza dentro del GraphQL Gateway} (tabla de la sección 8); que la ruta esté permitida no autoriza todas sus operaciones.

\subsection{Inventario simbólico de secretos}
Referencias, nunca valores. Todos los genera \texttt{vault-\allowbreak{}init} al azar (256 bits) y viven en Vault (\texttt{secret/\allowbreak{}}, KV v2), salvo \texttt{MONGO\_\allowbreak{}ROOT\_\allowbreak{}CLAVE}. Ninguno está en el repositorio, en el JavaScript del navegador, en capturas ni en logs. Cada servicio lee solo los suyos: su token de Vault tiene una política de solo lectura sobre esas rutas.

\vspace{0.6em}
{\scriptsize
\setlength{\LTpre}{0pt}\setlength{\LTpost}{0pt}
\setlength{\tabcolsep}{3pt}\renewcommand{\arraystretch}{1.15}
\begin{longtable}{|>{\raggedright\arraybackslash}p{\dimexpr 0.1393\linewidth-2\tabcolsep\relax}|>{\raggedright\arraybackslash}p{\dimexpr 0.0928\linewidth-2\tabcolsep\relax}|>{\raggedright\arraybackslash}p{\dimexpr 0.1184\linewidth-2\tabcolsep\relax}|>{\raggedright\arraybackslash}p{\dimexpr 0.0926\linewidth-2\tabcolsep\relax}|>{\raggedright\arraybackslash}p{\dimexpr 0.1295\linewidth-2\tabcolsep\relax}|>{\raggedright\arraybackslash}p{\dimexpr 0.1483\linewidth-2\tabcolsep\relax}|>{\raggedright\arraybackslash}p{\dimexpr 0.1309\linewidth-2\tabcolsep\relax}|>{\raggedright\arraybackslash}p{\dimexpr 0.1483\linewidth-2\tabcolsep\relax}|}
\hline
\textbf{Referencia} & \textbf{Tipo} & \textbf{Emisor} & \textbf{Consumidor (política)} & \textbf{Finalidad} & \textbf{Ubicación segura} & \textbf{Alcance} & \textbf{Rotación y revocación} \\ \hline
\endfirsthead
\hline
\textbf{Referencia} & \textbf{Tipo} & \textbf{Emisor} & \textbf{Consumidor (política)} & \textbf{Finalidad} & \textbf{Ubicación segura} & \textbf{Alcance} & \textbf{Rotación y revocación} \\ \hline
\endhead
\texttt{poke-\allowbreak{}fresh/\allowbreak{}internos/\allowbreak{}inventario} (\texttt{credencial}) & Credencial de servicio & \texttt{vault-\allowbreak{}init} & gateway, inventario, pedidos, graphql & Llamar a inventario & Vault & Solo inventario & Borrar el secreto y volver a correr \texttt{vault-\allowbreak{}init}; los servicios lo leen en $\leq$ 30~s \\ \hline
\texttt{poke-\allowbreak{}fresh/\allowbreak{}internos/\allowbreak{}clientes} (\texttt{credencial}) & Credencial de servicio & \texttt{vault-\allowbreak{}init} & gateway, clientes, pedidos, graphql & Llamar a clientes & Vault & Solo clientes & Igual \\ \hline
\texttt{poke-\allowbreak{}fresh/\allowbreak{}internos/\allowbreak{}pedidos} (\texttt{credencial}) & Credencial de servicio & \texttt{vault-\allowbreak{}init} & gateway, pedidos, pagos, graphql & Llamar a pedidos & Vault & Solo pedidos & Igual \\ \hline
\texttt{poke-\allowbreak{}fresh/\allowbreak{}internos/\allowbreak{}pagos} (\texttt{credencial}) & Credencial de servicio & \texttt{vault-\allowbreak{}init} & gateway, pagos, pedidos, graphql & Llamar a pagos & Vault & Solo pagos & Igual \\ \hline
\texttt{poke-\allowbreak{}fresh/\allowbreak{}internos/\allowbreak{}graphql} (\texttt{credencial}) & Credencial de servicio & \texttt{vault-\allowbreak{}init} & gateway, graphql & Llamar al GraphQL Gateway & Vault & Solo graphql & Igual \\ \hline
\texttt{poke-\allowbreak{}fresh/\allowbreak{}internos/\allowbreak{}auth-\allowbreak{}gateway} (\texttt{secreto}) & Secreto compartido (Basic) & \texttt{vault-\allowbreak{}init} & gateway, auth & El gateway ante el Auth Service & Vault & Rutas de sesión y usuarios del Auth Service & Rotar: el gateway lo relee en $\leq$ 30~s \\ \hline
\texttt{poke-\allowbreak{}fresh/\allowbreak{}internos/\allowbreak{}auth-\allowbreak{}clientes} (\texttt{secreto}) & Secreto compartido (Basic) & \texttt{vault-\allowbreak{}init} & clientes, auth & El servicio de clientes ante el Auth Service & Vault & Solo \texttt{/\allowbreak{}cuentas} (crear y cambiar cuentas de cliente) & Rotar \\ \hline
\texttt{poke-\allowbreak{}fresh/\allowbreak{}auth} (\texttt{pepper}) & Pepper de contraseñas & \texttt{vault-\allowbreak{}init} (una sola vez) & auth & Se mezcla con la clave antes de scrypt & Vault & Hashes de contraseñas & \textbf{No se rota a la ligera}: obliga a pedir cambio de contraseña a todos. Por eso \texttt{vault-\allowbreak{}init} nunca lo regenera \\ \hline
\texttt{poke-\allowbreak{}fresh/\allowbreak{}pasarela/\allowbreak{}comercio} (\texttt{api\_\allowbreak{}key}, \texttt{secreto\_\allowbreak{}firma}) & Claves del comercio ante el proveedor & \texttt{vault-\allowbreak{}init} en la demostración; el proveedor en producción & pagos, pasarela & Crear transacciones y firmar o verificar notificaciones & Vault & Operar la cuenta del comercio (nunca llega al navegador) & Pedir claves nuevas al proveedor y actualizar el secreto \\ \hline
\texttt{poke-\allowbreak{}fresh/\allowbreak{}mongo/\allowbreak{}<servicio>} (\texttt{usuario}, \texttt{clave}) & Credencial de base de datos & \texttt{vault-\allowbreak{}init} & ese servicio (y \texttt{semilla} la suya) & Conectarse a su base & Vault & \texttt{readWrite} solo en su base & Cambiar la clave en Vault y correr \texttt{vault-\allowbreak{}init} (actualiza el usuario en Mongo) \\ \hline
Token de Vault de cada servicio & Token con TTL 720 h & Vault, vía \texttt{vault-\allowbreak{}init} & ese servicio & Leer sus rutas & Volumen de solo lectura (ver «Distribución de tokens de Vault» en las brechas de la sección 0) & Solo lectura de sus rutas & Revocarlo en Vault; \texttt{vault-\allowbreak{}init} emite uno nuevo \\ \hline
\texttt{MONGO\_\allowbreak{}ROOT\_\allowbreak{}CLAVE} & Secreto de arranque & La persona que opera & mongo, vault-init & Crear los usuarios por servicio & \texttt{.\allowbreak{}env} fuera del repo & Administrar Mongo (ningún servicio la usa) & Cambiarla en Mongo y en \texttt{.\allowbreak{}env} \\ \hline
Certificados TLS y CA de prueba & Certificados & \texttt{tls-\allowbreak{}init} & todos los servicios & Cifrar y autenticar cada salto & Volúmenes \texttt{tls} (certificados) y \texttt{tls-\allowbreak{}ca} (clave de la CA, solo \texttt{tls-\allowbreak{}init}) & Un certificado por componente & Borrar el volumen \texttt{tls} y recrear: se emiten nuevos \\ \hline
\end{longtable}
}
\end{landscape}

\textbf{Si una clave se filtra:} se revoca o borra el secreto afectado y se corre \texttt{vault-\allowbreak{}init}; solo se ve comprometido lo que esa referencia permite (la columna Alcance), porque no hay un secreto común a todos los sistemas. Una credencial de pagos solo sirve para pagos; la de base de datos de un servicio no abre las de otro; el navegador nunca recibe ninguna.


\textbf{Si Vault no responde:} cada servicio sigue usando lo leído durante 5 minutos (\texttt{VAULT\_\allowbreak{}GRACIA\_\allowbreak{}SEGUNDOS}) y después responde 503 (falla cerrada). El gateway responde 503 \texttt{SECRETOS\_\allowbreak{}NO\_\allowbreak{}DISPONIBLES}. Probado en el caso 10.7.


\subsection{Identidades}

{\small
\setlength{\tabcolsep}{3pt}\renewcommand{\arraystretch}{1.15}
\begin{longtable}{|>{\raggedright\arraybackslash}p{\dimexpr 0.2485\linewidth-2\tabcolsep\relax}|>{\raggedright\arraybackslash}p{\dimexpr 0.4191\linewidth-2\tabcolsep\relax}|>{\raggedright\arraybackslash}p{\dimexpr 0.3323\linewidth-2\tabcolsep\relax}|}
\hline
\textbf{Identidad} & \textbf{Cómo se prueba} & \textbf{Dónde se valida} \\ \hline
\endfirsthead
\hline
\textbf{Identidad} & \textbf{Cómo se prueba} & \textbf{Dónde se valida} \\ \hline
\endhead
Usuario (cliente o personal) & Cookie con token opaco & Auth Service (introspección), por el gateway \\ \hline
Gateway ante servicios & \texttt{X-\allowbreak{}Credencial-\allowbreak{}Interna} del destino & Cada servicio, contra Vault \\ \hline
Gateway / clientes ante Auth & HTTP Basic (sobre HTTPS) con su secreto & Auth Service, contra Vault, y solo sus rutas \\ \hline
Servicio ante otro servicio & Credencial del destino + identidad del usuario original & El destino \\ \hline
Pasarela ante pagos & Firma HMAC del cuerpo & pagos \\ \hline
Pagos ante la pasarela & \texttt{X-\allowbreak{}Api-\allowbreak{}Key} & pasarela \\ \hline
Servicio ante Vault & Token con política propia & Vault \\ \hline
Servicio ante Mongo & Usuario y clave propios & Mongo (\texttt{authSource} = su base) \\ \hline
\end{longtable}
}

\subsection{Operaciones GraphQL}

{\footnotesize
\setlength{\tabcolsep}{3pt}\renewcommand{\arraystretch}{1.15}
\begin{longtable}{|>{\raggedright\arraybackslash}p{\dimexpr 0.2191\linewidth-2\tabcolsep\relax}|>{\raggedright\arraybackslash}p{\dimexpr 0.0956\linewidth-2\tabcolsep\relax}|>{\raggedright\arraybackslash}p{\dimexpr 0.1674\linewidth-2\tabcolsep\relax}|>{\raggedright\arraybackslash}p{\dimexpr 0.3188\linewidth-2\tabcolsep\relax}|>{\raggedright\arraybackslash}p{\dimexpr 0.1991\linewidth-2\tabcolsep\relax}|}
\hline
\textbf{Operación} & \textbf{Tipo} & \textbf{Quién} & \textbf{Servicios que combina} & \textbf{Pantalla} \\ \hline
\endfirsthead
\hline
\textbf{Operación} & \textbf{Tipo} & \textbf{Quién} & \textbf{Servicios que combina} & \textbf{Pantalla} \\ \hline
\endhead
\texttt{resumenPedido(\allowbreak{}id)} & query & dueño del pedido y personal & pedidos, pagos, inventario (+ clientes para admin) & \texttt{\#/\allowbreak{}orden/\allowbreak{}:\allowbreak{}id} \\ \hline
\texttt{misPedidos(\allowbreak{}primeros,\allowbreak{} despuesDe)} & query & cliente & pedidos & (disponible; la lista usa REST) \\ \hline
\texttt{reporteVentas(\allowbreak{}desde,\allowbreak{} hasta)} & query & admin, dueño & pedidos, inventario, pagos & \texttt{\#/\allowbreak{}admin/\allowbreak{}reportes} \\ \hline
\end{longtable}
}

No hay mutaciones en GraphQL: los cambios van por REST al servicio dueño, que aplica las reglas transaccionales.


\subsection{Métodos gRPC}

No aplica: no se implementó gRPC. La justificación, con los argumentos a favor y en contra y la evaluación de cada flujo, está en la sección 9 (“gRPC: por qué sí y por qué no”) y en el ADR 0001 (\texttt{docs/\allowbreak{}adr/\allowbreak{}0001-\allowbreak{}analitica-\allowbreak{}y-\allowbreak{}grpc.\allowbreak{}md}).

